% Subproblem 2: Rewrite the target integral as a quarter-turn defect
%
% CONTEXT: Part of the proof that Hess_p log J_8 > 0 on T_{2g}.
% Working directory: /data/projects/3eacb340-027e-4b71-8de4-67574ba54ea2/QBLYHzbdEjWeyg/proof/
%
% GIVEN (from Subproblem 1):
%   - psi = phi - pi/4; in (theta, psi) coordinates, x=sin(theta)cos(psi), y=sin(theta)sin(psi).
%   - F_A(theta,psi) = partial_psi G_A(theta,psi) where G_A = sum_{q in A} G_q.
%   - F_B(theta,psi) = F_A(theta, psi+pi).
%   - C(theta,psi) = F_A(theta,psi) + F_A(theta,psi+pi) = F_A + F_B.
%   - D(theta,psi) = F_A(theta,psi) - F_A(theta,psi+pi) = F_A - F_B.
%   - The full cube has C_4 rotational symmetry about the z-axis, i.e., it is invariant
%     under the map R_{pi/2}: psi -> psi + pi/2.
%   - u(theta,psi) = 2*pi*(G_A + G_B)(theta,psi) = 2*pi * sum_{all 8 cube vertices} G_q,
%     which satisfies u(theta, psi + pi/2) = u(theta, psi) (C_4 symmetry of the full cube).
%   - mu^pm are probability measures with densities rho^pm = exp(+/- u) / I^pm on S^2.
%   - rho(theta,psi) := rho^+(theta,psi) + rho^-(theta,psi) (not normalized, but positive).
%
% STATEMENT TO PROVE:
%
% (a) Show that rho(theta, psi + pi/2) = rho(theta, psi) for all (theta, psi).
%     (This follows from the C_4 symmetry u(theta, psi+pi/2) = u(theta, psi), which
%      implies exp(+/-u(psi+pi/2)) = exp(+/-u(psi)), and I^pm are constants.)
%
% (b) Show that C(theta, psi + pi/2) = C(theta, psi).
%     (This is a symmetry property of C under the quarter-turn.)
%
% (c) Define
%       Xi(theta, psi) = (1/2)*D(theta,psi)^2 + (1/2)*D(theta, psi+pi/2)^2 - C(theta,psi)^2.
%     Prove that the target integral
%       I = integral_{S^2} F_A(theta,psi) * F_A(theta,psi+pi) * d(mu^+ + mu^-)
%     satisfies
%       I = - integral_0^pi integral_0^{pi/2} Xi(theta,psi) * rho(theta,psi) * sin(theta) d psi d theta.
%
% PROOF STRATEGY for (c):
%   Use F_A * F_B = F_A * F_A(psi+pi) and the identities:
%     F_A = (C + D)/2,   F_B = F_A(psi+pi) = (C - D)/2.
%   So F_A * F_B = (C^2 - D^2)/4.
%
%   Partition the psi circle [0, 2*pi) into four quarter-turn sectors [0,pi/2), [pi/2,pi),
%   [pi,3*pi/2), [3*pi/2,2*pi). Use the pi/2-periodicity of rho and the transformation
%   behavior of C and D under psi -> psi + pi/2 to fold the integral over [0,2*pi)
%   down to [0,pi/2).
%
%   Key: Compute how D transforms under psi -> psi + pi/2.
%   Show that C(psi+pi/2) = C(psi) and D(psi+pi/2) can be expressed, then show
%   integral_0^{2*pi} F_A * F_B * rho dpsi = integral_0^{pi/2} (something involving C,D) * rho dpsi.
%   The formula Xi = (1/2)D^2 + (1/2)(D(psi+pi/2))^2 - C^2 should emerge naturally.
%
% Note: The identity F_A*F_B = (C^2 - D^2)/4 gives:
%   4 * integral F_A*F_B * rho dpsi = integral (C^2 - D^2) * rho dpsi.
%
%   Over [0,2*pi], using pi/2-periodicity of rho and C:
%   integral_0^{2*pi} C^2 * rho dpsi = 4 * integral_0^{pi/2} C^2 * rho dpsi.
%   integral_0^{2*pi} D^2 * rho dpsi = integral_0^{pi/2} [D^2 + D(psi+pi/2)^2 + D(psi+pi)^2 + D(psi+3pi/2)^2] * rho dpsi.
%   Use D(psi+pi) = -D(psi) (anti-symmetry) and D(psi+pi+pi/2) = -D(psi+pi/2) to get:
%   integral_0^{2*pi} D^2 * rho dpsi = 2 * integral_0^{pi/2} [D^2 + D(psi+pi/2)^2] * rho dpsi.
%
% DELIVERABLE: A clean LaTeX proof of (a)-(c). The output should be a self-contained
% LaTeX section (no \documentclass) ready for \input{} into the main paper.
\end{document}
